% TRANSLATION-PRODUCER DRAFT — UNCHECKED
% Work: Emmy Noether, Paper 41 — Korean translation unit U06
% Source custody snapshot: C:\Users\Floris\Documents\interlanguage\03_projects\language_management\cjk\03_working_translations\noether_paper41_zh_translation_001_20260722\source\Noether_Paper41_CurrentGermanAuthority_interval.tex
% Source snapshot lines: 68--78 (historical whole-source lines 19776--19786)
% Source unit: 2,642 LF-normalized UTF-8 bytes; SHA-256 1EC47AC260130703D0C0313186283365D339F04472184000D219ED44962EA55D
% Source snapshot SHA-256: C265058425E5E2D1A2289CC03A9DDEDDDF4803A3215DC3F173B93E7AB69D60ED
% Historical whole-source SHA-256: 443EF950D7D45DC6D9E44A9B87501620C10DA873E50E5F2B253ECCAE6A946D27
% Producer role: translation only. No source/scan check, Korean review, compilation, rendering, assembly, packaging, certification, approval, or self-check was performed.
% Handoff state: UNCHECKED; requires an independent Korean checker.

먼저 제1정식을 증명한 다음 그것이 제2정식 및 제3정식과 일치함을 보이겠다. 이 증명은 주어진 종류의 군 자기동형사상이 모두 \(A\)의 환 자기동형사상으로 연장될 수 있고, 이 환 자기동형사상이 알려진 바와 같이 내부자기동형사상이라는 사실에 근거한다. 실제로 이 자기동형사상 아래에서 \(u_S\)의 상을 \(v_S\)라 하자. 그러면 가정에 따라 관계 (2)부터 (4)가 모두 보존되므로, \(\sum v_SK\)도 같은 인자계에 대한 \(\mathfrak G\)와 \(K\)의 교차곱을 이룬다. 대응 \(\sum u_Sb_S\mapsto\sum v_Sb_S\)은 \(K\)의 원소들을 자기 자신에 대응시키는 환 자기동형사상을 준다. 따라서 앞에서 \(\mathfrak G^*\)에 관하여 말한 바에 의해, 이 자기동형사상은 \(K^*\)의 한 원소에 의해 유도된다.

제2정식으로 넘어가는 근거는 (2)에 의해 \(v_S\)가 \(v_S=u_Sc_S\) 꼴이어야 한다는 데 있다. 또한 인자계도 보존되므로 (5)에 따라 그에 딸린 변환량들은 모두 1이어야 한다. 역으로 이것이 성립하면, (2)부터 (5)는 다시 치환 \(v_S=u_Sc_S\)가 앞에서 지정한 종류의 \(\mathfrak G^*\)의 자기동형사상을 준다는 것을 보인다. 따라서 제1정식에 의해
\[
v_S=b^{-1}u_Sb=u_Sb/b^S\qquad\text{이고 따라서}\qquad c_S=b^{1-S}.
\]
특히 순환인 경우에는 제2정식의 가정이 (5')에 따라 \(N(c)=1\)로 표현되며, \(c=b^{1-S}\)는 잘 알려진 정리를 준다\srcfn{9)}{여기서 제시한 증명은 특히 \(k\)가 유한개의 원소만을 가질 때에도 성립한다. 이 경우 힐베르트의 증명은 성립하지 않지만 슈파이저의 증명은 여전히 유효하다.}.

제3정식은 제2정식과 곧바로 동치이다. 실제로 제2정식의 가정은 \(u_S\mapsto c_S\)가 인자계 1에 대한 1차 교차표현을 이룬다는 뜻이다. 명제 \(c_S=b^{1-S}\)는 이 표현이 일표현과 동치임, 즉 \(c_S\sim1\)임을 뜻한다.

다만 이 제3정식은 훨씬 더 일반적인 다음 정리의 특수한 경우일 뿐임을 지적해 둔다. 주어진 인자계 \(a_{S,T}\)에 대하여 불가약 교차표현류는 단 하나만 존재한다. 그 차수는 \(m\)이며, 여기서 \(m\)은 대응하는 대수의 슈어 지표\srcfn{10)}{I. Schur, Einige Bemerkungen zu der vorstehenden Arbeit des Herrn A. Speiser\textsuperscript{3)}. Math. Zeitschr. 5 (1919), S. 7--10. 6)에서 인용한 나의 강의에 나오는 교차표현 가군에 근거한 증명에 관해서는, 6)에서 인용한 M. Deuring의 보고서를 보라.}, 곧 대응하는 나눗셈대수의 차수이다.
